\section{Noetherianity and finiteness}\label{HNsec}
In this section, we prove Theorem~\ref{HN}. Recall that by Noetherian we mean right and left Noetherian. We begin with

\begin{teo} \label{LonNoeth} The algebras ${\mathcal L}_{0,n}$, $\Ll_{0,n}^A$ and $\Ll_{0,n}^\ep$, $\ep\in \mc^\times$, are Noetherian.
\end{teo}

By Proposition \ref{L0NAfgfree}, each of the algebras in this theorem is finitely generated.

Theorem \ref{LonNoeth} for ${\mathcal L}_{0,1}$ and any $\mathfrak{g}$ follows immediately from Joseph--Letzter's Theorem \ref{JLteo1}, claim (3), by identifying ${\mathcal L}_{0,1}$ with \smash{$U_q^{\rm lf}$} via $\Phi_1$. The method of proof uses filtration arguments. An alternative proof in the case of $\mathfrak{sl}(n)$, which works also for ${\mathcal L}_{0,1}^A$, was obtained by Domokos--Lenagan in \cite{DomLen}, by exhibiting special sequences of generators of \smash{${\mathcal L}_{0,1}^A$} satisfying {\it polynormal} relations, as we define now.

\begin{defi}[{see \cite[Proposition 3.133]{VY}}] Let $R$ be a Noetherian Abelian ring, and~$B$ a~finitely generated $R$-algebra with product $\circ$. We call polynormal a set of relations between generators $u_1,\dots,u_M$ of~$B$, of the form{\samepage
\begin{equation}\label{echangegrad++}
u_i\circ u_j - q_{ij} u_j\circ u_i = \sum_{s=1}^{j-1}\sum_{t=1}^M \bigl(\alpha_{ij}^{st} u_s\circ u_t + \beta_{ij}^{st} u_t\circ u_s\bigr)
\end{equation}
for all $1\leq j<i\leq M$, where $\alpha_{ij}^{st}, \beta_{ij}^{st} \in R$, and the elements $q_{ij}\in R$ are invertible.}
\end{defi}
Note that this definition is more restrictive than the more standard one, e.g., in~\cite[Definition~II.4.1]{BG}. If such a set of relations exists in $B$, then $B$ can be endowed with an algebra filtration such that the associated graded algebra is a quotient of a skew-polynomial algebra~\cite[Proposition I.8.17]{BG}. By classical results, we have (see, e.g., \cite[Theorems~1.2.9, 1.6.9 and Examples~1.6.11]{MC-R}, or~\cite[Lemmas~3.130--3.131]{VY}):

\begin{teo} If the algebra filtration is well founded, then $B$ is a Noetherian ring.
\end{teo}

In \cite{DomLen}, Theorem \ref{LonNoeth} is also proved for any $n\geq 1$ in the case of $\mathfrak{g}={\mathfrak{sl}_2}$ by considering~$\Ll_{0,n}^A({\mathfrak{sl}_2})$ as an iterated overring of ${\mathcal L}_{0,1}({\mathfrak{sl}_2})$.


The proof of Theorem \ref{LonNoeth} that we develop for any $\mathfrak{g}$ and $n\geq 1$ is also based on polynormal relations. In our proof, the generating set of $\Ll_{0,n}$ that we will consider is evident, as they are matrix coefficients in the modules $V_{\varpi_k}$, $k\in\{1,\dots,m\}$; the task is then to exhibit a set of polynormal relations between them, that hold in a certain graded algebra associated to $\Ll_{0,n}$. Indeed, as explained above this will imply that the graded algebra is Noetherian, and that $\Ll_{0,n}$ is Noetherian as well. In the case of $\Ll_{0,n}^A$, the proof is formally similar, but it needs the use of canonical bases discussed in Section \ref{canbasemodqg}.

\begin{proof}[Proof of Theorem \ref{LonNoeth}] First, we develop the proof for $\Ll_{0,n}$, and then for $\Ll_{0,n}^A$; the result for%
\[\Ll_{0,n}^\ep\!= \Ll_{0,n}^A/(q-\ep)\Ll_{0,n}^A
\]
 follows immediately by lifting ideals by the quotient map \smash{$\Ll_{0,n}^A \!\ra \Ll_{0,n}^\ep$}.

We adapt the proof of Theorem \ref{JLteo1}\,(3) given in \cite[Theorem 3.137]{VY}. Let us begin by recalling these arguments. In doing this, let us stress that \cite{VY} takes on ${\mathcal O}_q$ and $\Ll_{0,1}$ the product opposite to ours, and below in \eqref{echangegrad} and \eqref{echangegrad+} we respect their convention.

As usual, let $C(\mu)$ be the vector space generated by the matrix coefficients of $V_\mu$, the simple~$U_q^{\rm ad}$-module of highest weight $\mu \in P_+$. Denote by $C(\mu)_\lambda \subset C(\mu)$ the subspace of weight $\lambda$ for the left coregular action of $U_q({\mathfrak h})$; so $\alpha\in C(\mu)_\lambda$ if \smash{$K_\nu\rhd \alpha = q^{(\nu,\lambda)}\alpha$}, $\nu\in P$. Consider the semigroup
\[\Lambda=\{(\mu, \lambda)\in P_+\times P, \, \lambda \;\text{is a weight of}\;V_\mu\}.\]
Recall that the partial order $\preceq$ on $P$ is defined by $\mu\preceq \mu'$ if and only if $\mu'-\mu\in D^{-1}Q_+$. Define~$\preceq$ on $\Lambda$ by: $(\mu,\lambda)\preceq (\mu', \lambda')$ if and only if $\mu'-\mu\in D^{-1}Q_+$ and $\lambda'-\lambda\in D^{-1}Q_+$. If~${(\mu,\lambda)\preceq (\mu', \lambda')}$ and $(\mu,\lambda)\ne (\mu', \lambda')$, we write $(\mu,\lambda)\prec (\mu', \lambda')$. Since ${\mathcal L}_{0,1}$ and ${\mathcal O}_q$ are isomorphic vector spaces, we have $ {\mathcal L}_{0,1}=\bigoplus_{\mu\in P_+} C(\mu)=\bigoplus_{(\mu,\lambda)\in \Lambda} C(\mu)_\lambda$. Consider the family of subspaces
\begin{align*}
& {\mathcal F}_2^{\mu,\lambda} :=\bigoplus_{(\mu',\lambda')\preceq (\mu,\lambda)} C(\mu')_{\lambda'} , \qquad{\mathcal F}_2^{\prec \mu,\lambda}:=\bigoplus_{(\mu',\lambda')\prec (\mu,\lambda)} C(\mu')_{\lambda'} , \qquad (\mu,\lambda)\in \Lambda.
\end{align*}
We have
\begin{equation}\label{filtrationMatthieu}
{\mathcal L}_{0,1} = \bigcup_{(\mu,\lambda)\in \Lambda} {\mathcal F}_2^{\mu,\lambda}.
\end{equation}
Indeed, clearly
\[ {\mathcal L}_{0,1} = \sum_{(\mu,\lambda)\in \Lambda} {\mathcal F}_2^{\mu,\lambda},
\]
 so \eqref{filtrationMatthieu} follows from the following fact: for every $(\mu,\lambda), (\mu',\lambda')\in \Lambda$, the element $(\mu'',\lambda''):= (\mu+\mu',\lambda+\lambda')$ is such that
\[{\mathcal F}_2^{\mu,\lambda}+ {\mathcal F}_2^{\mu',\lambda'} \subset {\mathcal F}_2^{\mu'',\lambda''}.\]
Note that in general, since $Q_+\nsubseteq P_+$ (but $P_+\subset D^{-1}Q_+$), it is not true that there exists an element $(\mu'',\lambda'')$ satisfying such an inclusion if one replaces $\preceq$ with the standard ``product'' partial order $\leq$ on $\Lambda$, defined by $(\mu,\lambda)\leq (\mu', \lambda')$ if and only if $\mu'-\mu\in Q_+$ and $\lambda'-\lambda\in Q_+$. Note also that $\preceq$ is finer than $\leq$, in the sense that if $\mu \leq \mu'$, then $\mu\preceq \mu'$. Again, this would not be true if we had replaced $D^{-1}Q_+$ by $P_+$ in the definition of $\preceq$.

The family \smash{${\mathcal F}_2:=\big\{{\mathcal F}_2^{\mu,\lambda}\big\}_{(\mu,\lambda)\in \Lambda}$} is a filtration of the vector space $ {\mathcal L}_{0,1}$, which is clearly well founded (i.e., every subset of $\Lambda$ contains a minimal element, or equivalently any decreasing infinite sequence of elements in $\Lambda$ is eventually constant).

Consider the associated graded vector space \smash{${\rm Gr}_{\mathcal F_2}({\mathcal L}_{0,1}) := \bigoplus_{(\mu,\lambda)} {\mathcal F}_2^{\mu,\lambda}/{\mathcal F}_2^{\prec \mu,\lambda}$}. By identifying an element $x\in C(\mu)_\lambda$ with its coset \smash{$\bar{x}\in {\mathcal F}_2^{\mu,\lambda}/{\mathcal F}_2^{\prec \mu,\lambda}$}, we get an equality of vector spaces $ {\rm Gr}_{{\mathcal F}_2}({\mathcal L}_{0,1})= \bigoplus_{(\mu,\lambda)\in \Lambda} C(\mu)_\lambda$. Now, one has the following facts:

(i) Taking the product in $\Ll_{0,1}$, we have
\begin{equation}\label{gradalg01}
\alpha \beta\in {\mathcal F}_2^{\mu_1+\mu_2,\lambda_1+\lambda_2}\qquad \mathrm{for}\quad \alpha \in C(\mu_1)_{\lambda_1},\quad \beta \in C(\mu_2)_{\lambda_2}.
\end{equation}
This follows from \eqref{decomprep} and the fact that, for every $\nu\in P_+$ and every summand of the formula~\eqref{mmtilde}, denoting by $-r\in -Q_+$ the weight of the $R$-matrix component $R_{(2)}$ we have
\begin{gather*}%\label{weightrhd}
K_\nu\rhd \bigl((R_{(2')}{S}(R_{(2)}) \rhd \alpha) \star (R_{(1')}\rhd \beta \lhd R_{(1)})\bigr) \\
\qquad= q^{(\nu,\lambda_1+\lambda_2-r)}(R_{(2')}{S}(R_{(2)}) \rhd \alpha) \star (R_{(1')}\rhd \beta \lhd R_{(1)}).
\end{gather*}
(Details of a similar computation are given below \eqref{prodL0nform-duplicate}.) It follows from~\eqref{gradalg01} that~$\mathcal{F}_2$ is an algebra filtration of $\Ll_{0,1}$, and then ${\rm Gr}_{\mathcal F_2}({\mathcal L}_{0,1})$ is a graded algebra.

(ii) Denote by $\alpha \circ \beta$ the product in $ {\rm Gr}_{{\mathcal F}_2}({\mathcal L}_{0,1})$ of $\alpha,\beta\in \Ll_{0,1}$. The space $C(\mu_1+\mu_2)$ has multiplicity one in $C(\mu_1)\otimes C(\mu_2)$ (again by~\eqref{decomprep}), therefore if $\alpha \in C(\mu_1)_{\lambda_1}$ and $\beta \in C(\mu_2)_{\lambda_2}$, then $\alpha \circ \beta$ is the projection of $\alpha \beta$ onto $C(\mu_1+\mu_2)_{\lambda_1+\lambda_2}$.
 Denote by $\bar{\star}$ the product $\star$ of $\Oo_q$ followed by the projection onto the component $C(\mu + \nu)$. Then, we have
\begin{equation}\label{filtration=}
C(\mu) \circ C(\nu) = C(\mu)\ \bar{\star} \ C(\nu) = C(\mu + \nu).
\end{equation}
This follows from the formula \eqref{mmtilde}, and the fact that it is given by an invertible twist of the product $\star$.

(iii) For every $\mu\in P_+$, fix a basis of weight vectors \smash{$e_1^\mu,\dots,e_{d(\mu)}^\mu$} of $V_\mu$. Denote by $e^1_\mu,\dots,e^{d(\mu)}_\mu\!\in V_\mu^*$ the dual basis, and by $w\big(e_i^\mu\big)$ the weight of $e_i^\mu$. Consider the matrix coefficients ${}_{\mu}\phi^i_j(x) :=\smash{ e^i_\mu\big(\pi_V(x)\big(e_j^\mu\big)\big)}$, $x\in U_q$. By using the formula \eqref{mmtilde} and the explicit form of the $R$-matrix, one can check that
\begin{align}
{}_{{\mu}}\phi^i_j \circ {}_{{\nu}}\phi^k_l & = \sideset{}{'}\sum_{j',l'} c_{j',l'}^{ikjl}\ {}_{{\mu}}\phi^i_{j'} \ \bar{\star} \ {}_{{\nu}}\phi^k_{l'}\nonumber \\ & = q^{(w(e_j^\mu),w(e_l^\nu)-w(e_k^\nu))} {}_{{\mu}}\phi^i_{j} \ \bar{\star} \ {}_{{\nu}}\phi^k_{l} + \sideset{}{'}\sum_{\substack{j',l'\\ j'\ne j,\, l'\ne l}} d_{j',l'}^{ikjl}\ {}_{{\mu}}\phi^i_{j'} \circ {}_{{\nu}}\phi^k_{l'},\label{prodcirc1}
\end{align}
where \smash{$\sum^{'}_{j',l'}$} is the sum over indices with weights satisfying
\[
w\big(e_j^\mu\big) + w\big(e_l^\nu\big) = w\big(e_{j'}^\mu\big) + w\big(e_{l'}^\nu\big), \qquad w\big(e_{j'}^\mu\big)\leq w\big(e_{j}^\mu\big)
\qquad \text{and} \qquad w\big(e_{l'}^\nu\big)\geq w\big(e_{l}^\nu\big),
\]
 and the coefficient \smash{$c_{j,l}^{ikjl}$}, equal to \smash{$q^{(w(e_j^\mu),w(e_l^\nu)-w(e_k^\nu))}$}, is computed from the term $\Theta$ in the $R$-matrix factorization \eqref{Rmatfact}. In general, all the coefficients~\smash{$c_{j',l'}^{ikjl}$} and \smash{$d_{j',l'}^{ikjl}$} belong to $\mc(q)$ (see \cite[Proposition 4.1]{BR}); in particular \smash{$q^{(w(e_j^\mu),w(e_l^\nu)-w(e_k^\nu))}\in q^\mz$} since $w(e_l^\nu)-w(e_k^\nu)\in Q$. The second equality follows by repeated use of the first and \eqref{filtration=}. Similarly, by using \eqref{mmtildebis} one gets
\begin{align*}
{}_{{\nu}}\phi^k_l \circ {}_{{\mu}}\phi^i_j & = \sideset{}{'}\sum_{i',k'} e_{i',k'}^{kilj}\ {}_{{\mu}}\phi^{i'}_{j} \ \bar{\star} \ {}_{{\nu}}\phi^{k'}_{l}\\
 & = q^{(w(e_i^\mu),w(e_k^\nu)-w(e_l^\nu))} {}_{{\mu}}\phi^i_{j} \ \bar{\star} \ {}_{{\nu}}\phi^k_{l} + \sideset{}{'}\sum_{\substack{i',k'\\ i'\ne i,\, k'\ne k}} e_{i',k'}^{kilj}\ {}_{{\mu}}\phi^{i'}_{j} \ \bar{\star}\ {}_{{\nu}}\phi^{k'}_{l}\\
 & = q^{(w(e_i^\mu),w(e_k^\nu)-w(e_l^\nu))} {}_{{\mu}}\phi^i_{j} \ \bar{\star} \ {}_{{\nu}}\phi^k_{l} + \sideset{}{'}\sum_{\substack{i',k',j',l'\\ i'\ne i,\, k'\ne k}} f_{i',k'}^{kilj}\ {}_{{\mu}}\phi^{i'}_{j'} \ \circ\ {}_{{\nu}}\phi^{k'}_{l'},
\end{align*}
where \smash{$e_{i',k'}^{kilj}$}, \smash{$f_{i',k'}^{kilj}\in \mc(q)$}, and \smash{$\sum^{'}_{i',k'}$} is the sum over indices with weights satisfying \begin{gather*}
w\big(e_i^\mu\big) + w(e_k^\nu) = w\big(e_{i'}^\mu\big) + w(e_{k'}^\nu), \qquad w\big(e_{i'}^\mu\big)\leq w\big(e_{i}^\mu\big),\\
w(e_{k'}^\nu)\geq w(e_{k}^\nu) , \qquad e_{i,k}^{kilj} = q^{(w(e_i^\mu),w(e_k^\nu)-w(e_l^\nu))}.
\end{gather*}
 The third equality comes from the second and \eqref{prodcirc1}; the sum is over indices with weights satisfying
 \begin{gather*}
 w\big(e_i^\mu\big) + w(e_k^\nu) = w\big(e_{i'}^\mu\big) + w(e_{k'}^\nu), \\
 w\big(e_{i'}^\mu\big)< w\big(e_{i}^\mu\big), \qquad w(e_{k'}^\nu)> w(e_{k}^\nu), \qquad w\big(e_{j'}^\mu\big)\leq w\big(e_{j}^\mu\big), \qquad w(e_{l'}^\nu)\geq w(e_{l}^\nu).
 \end{gather*} By eliminating the leading term \smash{${}_{{\mu}}\phi^i_{j} \ \bar{\star} \ {}_{{\nu}}\phi^k_{l}$}, one deduces
\begin{align}\label{commutphiklij}
&{}_{{\nu}}\phi^k_l \circ {}_{{\mu}}\phi^i_j - q_{ijkl} \ {}_{{\mu}}\phi^i_j \circ {}_{{\nu}}\phi^k_l = \sideset{}{'}\sum_{\substack{i',k',j',l'\\ i'\ne i,\, k'\ne k}} f_{i',k'}^{kilj}\ {}_{{\mu}}\phi^{i'}_{j'} \ \circ\ {}_{{\nu}}\phi^{k'}_{l'}-\sideset{}{'}\sum_{\substack{j',l'\\ j'\ne j,\, l'\ne l}} q_{ijkl} d_{j',l'}^{ikjl}\ {}_{{\mu}}\phi^i_{j'} \circ {}_{{\nu}}\phi^k_{l'},\!\!\!
\end{align}
where $q_{ijkl} = q^{(w(e_j^\mu)+ w(e_i^\mu),w(e_k^\nu)-w(e_l^\nu))}$.

(iv) We can always reorder the weight vectors $e_1^\mu,\dots,e_{d(\mu)}^\mu$ so that $w\big(e_i^\mu\big) > w\big(e_j^\mu\big)$ implies $i < j$; then \eqref{commutphiklij} reads
\begin{align}
{}_{{\nu}}\phi^k_l \circ {}_{{\mu}}\phi^i_j - q_{ijkl} \ {}_{{\mu}}\phi^i_j\circ {}_{{\nu}}\phi^k_l ={}& \sum_{r=i}^{d(\mu)}\sum_{s=1}^k \sum_{u=1}^{l-1} \sum_{v=j+1}^{d(\mu)} \delta^{ijkl}_{rsuv} \ {}_{\mu}\phi^r_v \circ {}_{{\nu}}\phi^s_u \nonumber\\
 & - \sum_{r=i+1}^{d(\mu)}\sum_{s=1}^{k-1} q_{ijkl}\gamma_{rs}^{ijkl} \ {}_{\mu}\phi^r_j \circ {}_{{\nu}}\phi^s_l, \label{echangegrad}
\end{align}
where \smash{$\gamma_{rs}^{ijkl}, \delta^{ijkl}_{rsuv}\in \mc(q)$} are such that $\gamma_{rs}^{ijkl}=0$ unless $w(e_r^\mu) < w(e_i^\mu)$ and $w(e_s^\nu) > w(e_k^\nu)$, and \smash{$\delta^{ijkl}_{rsuv}=0$} unless $w(e_u^\nu) > w(e_l^\nu)$, $w(e_v^\mu) < w(e_j^\mu)$, $w(e_r^\mu) \leq w(e_i^\mu)$ and $w(e_s^\nu) \geq w(e_k^\nu)$. Now, from \eqref{echangegrad} one can extract a defining set of polynormal relations for $ {\rm Gr}_{{\mathcal F}_2}({\mathcal L}_{0,1})$, as in \eqref{echangegrad++}. Indeed, like ${\mathcal L}_{0,1}$ the algebra $ {\rm Gr}_{{\mathcal F}_2}({\mathcal L}_{0,1})$ is generated by the matrix coefficients ${}_{{\varpi_k}} \phi_i^j$ of the fundamental representations $V_{\varpi_k}$. One can list these matrix coefficients, say $M$ in number, in an ordered sequence $u_1,\dots,u_M$ such that the following condition holds: if $w(e_k^{\varpi_s}) < w(e_i^{\varpi_r})$, or~${w(e_k^{\varpi_s}) = w(e_i^{\varpi_r})}$ and $w(e_l^{\varpi_s}) < w(e_j^{\varpi_r})$, then $u_a:={}_{{\varpi_r}} \phi^i_j$ and $u_b:={}_{{\varpi_s}} \phi^k_l$ satisfy $b<a$. Then denoting \smash{${}_{{\mu}}\phi^i_j$}, \smash{${}_{{\nu}}\phi^k_l$} in \eqref{echangegrad} by $u_j$, $u_i$, respectively, and assuming $u_j<u_i$, one finds that all terms~${u_s:= {}_{\mu}\phi^r_v}$, ${}_{\mu}\phi^r_j$ in the sums are $< u_j$. Therefore, for all $1\leq j<i\leq M$ it takes the form
\begin{equation}\label{echangegrad+}
u_i\circ u_j - q_{ij} u_j\circ u_i = \sum_{s=1}^{j-1}\sum_{t=1}^M \alpha_{ij}^{st} u_s\circ u_t
\end{equation}
for some $q_{ij}\in q^{\mz}$ and $\alpha_{ij}^{st}\in \mc(q)$. As explained after \eqref{echangegrad++}, it follows that $ {\rm Gr}_{{\mathcal F}_2}({\mathcal L}_{0,1})$ is a~Noetherian ring, and since the filtration $\mathcal{F}_2$ is well founded, it implies that $\Ll_{0,1}$ is Noetherian too.

We are going to extend all these facts to $\Ll_{0,n}$, $n>1$. First, we need to refine the filtration~${\mathcal F}_2$ on $\Ll_{0,1}$. Consider the action of $U_q({\mathfrak h})$ on $C(\mu)_{\lambda}$ given by
\begin{equation}\label{coadweightnew}
K_\nu.\alpha := {\rm coad}\big(K_\nu^{-1}\big)(\alpha),\qquad \nu\in P, \quad \alpha \in C(\mu)_{\lambda}.
\end{equation}
Denote by $C(\mu)_{\lambda,\gamma} \subset C(\mu)_{\lambda}$ the subspace of weight $\gamma$ for
this action; so $\alpha \in C(\mu)_{\lambda,\gamma}$ if $K_\nu.\alpha= q^{(\nu,\gamma)}\alpha$.
Consider the semigroup
\[\Lambda_P=\{(\mu, \lambda,\gamma)\in P_+\times P^2, \, \lambda \;\text{is a weight of}\;V_\mu \;\text{for} \; \rhd ,\, \gamma \;\text{is a weight of}\;V_\mu \;\text{for}\ . \}\]
with the partial order $(\mu,\lambda,\gamma)\preceq (\mu', \lambda',\gamma')$ if and only if $\mu'-\mu,\lambda'-\lambda,\gamma'-\gamma\in D^{-1}Q_+$. Define
\begin{gather*}
[\Lambda_P] =\bigl\lbrace([\mu],[\lambda],[\gamma])\in P_+^n\times P^n\times P^n \\
\hphantom{[\Lambda_P] =\bigl\lbrace([\mu],[\lambda],[\gamma]}{} \mid (\mu_i,\lambda_i,\gamma_i)\in \Lambda_P,\ [\mu] = (\mu_i)_{i=1}^n, [\lambda] = (\lambda_i)_{i=1}^n, [\gamma]=(\gamma_i)_{i=1}^n\bigr\rbrace .
\end{gather*}
Let us put the following lexicographic order on $ [\Lambda_P]$, starting from the tail: $([\mu'],[\lambda'],[\gamma'])\preceq ([\mu], [\lambda],[\gamma])$ if $(\mu_n',\lambda_n',\gamma_n')\prec (\mu_n, \lambda_n,\gamma_n)$, or $(\mu_n,\lambda_n,\gamma_n)= (\mu'_n, \lambda'_n,\gamma'_n)$ and $(\mu_{n-1}',\lambda_{n-1}',\gamma_{n-1}')\prec (\mu_{n-1}, \lambda_{n-1},\gamma_{n-1}),\dots$, or $(\mu_k,\lambda_k,\gamma_k)= (\mu'_k, \lambda'_k,\gamma'_k)$ for all $1< k\leq n$ and $(\mu_1',\lambda_1',\gamma_1')\preceq (\mu_1, \lambda_1,\gamma_1)$. (As usual, we write $([\mu'],[\lambda'],[\gamma'])\prec ([\mu], [\lambda],[\gamma])$ for $([\mu'],[\lambda'],[\gamma'])\preceq ([\mu], [\lambda],[\gamma])$ and $([\mu'],[\lambda'],[\gamma'])\ne ([\mu], [\lambda],[\gamma])$.)

Now recall that $\Ll_{0,n} = {\mathcal L}_{0,1}^{\otimes n} = {\mathcal O}_q^{\otimes n}$ as vector spaces. For every $([\mu],[\lambda],[\gamma])\in[\Lambda_P]$, consider the subspace $C([\mu])_{[\lambda],[\gamma]}\subset \Ll_{0,n}$ defined by
\begin{align*}
&C([\mu]) =C(\mu_1)\otimes \dots \otimes C(\mu_n),\qquad
C([\mu])_{[\lambda],[\gamma]}=C(\mu_1)_{\lambda_1,\gamma_1}\otimes \dots \otimes C(\mu_n)_{\lambda_n,\gamma_n}.
\end{align*}
Then $ {\mathcal L}_{0,n}=\bigoplus_{[\mu]\in P_+^n} C({[\mu]})$ and $ C({[\mu]})=\bigoplus_{([\lambda],[\gamma])} C([\mu])_{[\lambda],[\gamma]}$. For every
$([\mu],[\lambda],[\gamma])\in [\Lambda_P]$ define
\begin{align}
&{\mathcal F}_3^{[\mu],[\lambda],[\gamma]} =\bigoplus_{([\mu'],[\lambda'],[\gamma'])\preceq ([\mu],[\lambda],[\gamma])} C([\mu'])_{[\lambda'],[\gamma']},\label{F2grad}\\
&{\mathcal F}_3^{\prec [\mu],[\lambda],[\gamma]} =\bigoplus_{([\mu'],[\lambda'],[\gamma'])\prec ([\mu],[\lambda],[\gamma])} C([\mu'])_{[\lambda'],[\gamma']}.\notag
\end{align}
Clearly, $\Ll_{0,n}$ is the union of the subspaces ${\mathcal F}_3^{[\mu],[\lambda],[\gamma]}$ over all $([\mu],[\lambda],[\gamma])\in [\Lambda_P]$, so these form a vector space filtration of $\Ll_{0,n}$. Let us denote it ${\mathcal F}_3$, and define
\[{\rm Gr}_{{\mathcal F}_3}({\mathcal L}_{0,n})_{[\mu],[\lambda],[\gamma]} = {\mathcal F}_3^{[\mu],[\lambda],[\gamma]}/{\mathcal F}_3^{\prec [\mu],[\lambda],[\gamma]}.\]
This space is canonically identified with $C({[\mu]})_{[\lambda],[\gamma]}$, so the graded vector space associated to~$\mathcal{F}_3$~is%
\begin{equation}\label{decompgrad}
{\rm Gr}_{{\mathcal F}_3}({\mathcal L}_{0,n}) = \bigoplus_{([\mu],[\lambda],[\gamma])\in [\Lambda_P]} {\rm Gr}_{{\mathcal F}_3}({\mathcal L}_{0,n})_{[\mu],[\lambda],[\gamma]} = \bigoplus_{([\mu],[\lambda],[\gamma])\in [\Lambda_P]} C({[\mu]})_{[\lambda],[\gamma]}.
\end{equation}
 We claim that ${\mathcal F}_3$ is an algebra filtration with respect to the product of $\Ll_{0,n}$, and therefore~${{\rm Gr}_{{\mathcal F}_3}({\mathcal L}_{0,n})}$ is a graded algebra.

For notational simplicity, let us prove it for $n=2$, the general case being strictly similar. Recall the $R$-matrix factorization \eqref{Rmatfact}. Take tuples $([\mu],[\lambda],[\gamma]) =((\mu_1,\mu_2),(\lambda_1,\lambda_2),(\gamma_1,\gamma_2))$ and $([\mu'],[\lambda'],[\gamma'])=((\mu_1',\mu_2'),(\lambda_1',\lambda_2'),(\gamma_1',\gamma_2'))$ in $[\Lambda_P]$, and elements $\alpha\otimes \beta \in C([\mu])_{[\lambda],[\gamma]}$ and~${\alpha'\otimes \beta'\in C([\mu'])_{[\lambda'],[\gamma']}}$. Recall from \eqref{prodL0nform} that the product of $\Ll_{0,2}$ is given by the formula
\begin{gather}
(\alpha\otimes \beta) (\alpha' \otimes \beta') \nonumber\\
\qquad= \sum_{(R^1),\dots, (R^4)} \alpha \bigl(S(R^3_{(1)}R^4_{(1)})\rhd \alpha' \lhd R^1_{(1)}R^2_{(1)}\bigr)\otimes \bigl( S(R^1_{(2)}R^3_{(2)}) \rhd \beta \lhd R^2_{(2)}R^4_{(2)}\bigr)\beta'.\label{prodL0nform-duplicate}
\end{gather}
For every $\nu\in P$ and any of the components \smash{$R^1_{(2)},\dots,R^4_{(2)}$}, denoting by $-r_j\in -Q_+$ the weight of $R^j_{(2)}$, we have
\begin{gather*}
K_\nu \rhd \big( S\big(R^1_{(2)}R^3_{(2)}\big) \rhd \beta \lhd R^2_{(2)}R^4_{(2)}\big) \\
 \qquad= \sum_{(\beta)} \beta_{(1)}\big(R^2_{(2)}R^4_{(2)}\big)\big(K_\nu S\big(R^1_{(2)}R^3_{(2)}\big)\rhd \beta_{(2)}\big)\\
 \qquad= q^{-(\nu,r_1+r_3)}\sum_{(\beta)} \beta_{(1)}\big(R^2_{(2)}R^4_{(2)}\big)\big(S\big(R^1_{(2)}R^3_{(2)}\big)K_\nu \rhd \beta_{(2)}\big)\\
 \qquad= q^{(\nu,\lambda_2-r_1-r_3)}\sum_{(\beta)} \beta_{(1)}\bigl(R^2_{(2)}R^4_{(2)}\bigr)\big(S\big(R^1_{(2)}R^3_{(2)}\big)\rhd \beta_{(2)}\big)\\
 \qquad= q^{(\nu,\lambda_2-r_1-r_3)} \big( S\big(R^1_{(2)}R^3_{(2)}\big) \rhd \beta \lhd R^2_{(2)}R^4_{(2)}\big).
\end{gather*}
By similar computations for the action ${\rm coad}(K_\nu^{-1})$, and for all terms in the right-hand side of \eqref{prodL0nform-duplicate}, and using \eqref{gradalg01} componentwisely, we find that
\[\alpha \big(S\big(R^3_{(1)}R^4_{(1)}\big)\rhd \alpha' \lhd R^1_{(1)}R^2_{(1)}\big)\otimes \big( S\big(R^1_{(2)}R^3_{(2)}\big) \rhd \beta \lhd R^2_{(2)}R^4_{(2)}\big)\beta' \in \Ff_3^{[\mu]+[\mu'], [\lambda''],[\gamma'']},\]
where
\begin{align*}
&\lambda'' = (\lambda_1+\lambda_1'+r_3+r_4,\lambda_2+\lambda_2'-r_1-r_3),\\
& \gamma'' = (\gamma_1+\gamma_1'+r_1+r_2+r_3+r_4,\gamma_2+\gamma_2'-r_1-r_2-r_3-r_4).
\end{align*}
Since $r_1+r_2+r_3+r_4=0$ implies $r_1=r_2=r_3=r_4=0$, by the order we have put on $[\Lambda_P]$, we deduce
 \[(\alpha\otimes \beta)(\alpha'\otimes \beta') \in \Ff_3^{[\mu]+[\mu'],[\lambda]+[\lambda'],[\gamma]+[\gamma']}.\]
Note that the filtration $\mathcal{F}_3$, taking the action \eqref{coadweightnew} into account, is crucial for this argument to work. Similar arguments work for any $n\geq 2$. This proves that ${\rm Gr}_{{\mathcal F}_3}({\mathcal L}_{0,n})$ is a graded algebra. We denote its product by $\circ_n$.

Next, we show that \eqref{filtration=} implies the analogous property for the product $\circ_n$. For simplicity of notations let us again assume that $n=2$. Recall that the product $\circ_2$ is defined on homogeneous elements \smash{$\overline{\alpha\otimes \beta} \in {\rm Gr}_{{\mathcal F}_3}({\mathcal L}_{0,n})_{[\mu],[\lambda]}$} and \smash{$\overline{\alpha'\otimes \beta'} \in {\rm Gr}_{{\mathcal F}_3}({\mathcal L}_{0,n})_{[\mu'],[\lambda']}$} by \[\overline{\alpha\otimes \beta}\circ_n\overline{\alpha'\otimes \beta'} = (\alpha\otimes \beta)(\alpha'\otimes \beta') + {\mathcal F}_3^{\prec [\mu+\mu'],[\lambda+\lambda']}.\]
Clearly, \eqref{filtration=} gives $(C(\mu_1)\circ C(\mu_1'))\otimes (C(\mu_2)\circ C(\mu_2')) = C([\mu + \mu'])$, and \eqref{prodL0nform-duplicate} gives
 \[C([\mu]) \circ_n C([\mu']) \subset (C(\mu_1)\circ C(\mu_1'))\otimes (C(\mu_2)\circ C(\mu_2')).\]
The converse inclusion holds true as well, as one can see by expressing, reciprocally, the (componentwise) product of $\Ll_{0,1}^{\otimes n}$ in terms of the product of $\Ll_{0,n}$ via the formula \eqref{prodL0nform3}. In conclusion,%
\[%\label{filtration=2}
C([\mu]) \circ_n C([\mu']) = C([\mu + \mu']).
\]
We are left to show that \eqref{echangegrad} generalizes to $\Ll_{0,n}$. First, note that for every $1\leq a\leq n$ the embedding $\mathfrak{i}_a\colon \Ll_{0,1}\ra \Ll_{0,n}$ in \eqref{plgtL01a} is a morphism of the filtered algebras $(\Ll_{0,1},\mathcal{F}_2)$ and $(\Ll_{0,n},\mathcal{F}_3)$, in the sense that
\[\mathfrak{i}_a\big({\mathcal F}_2^{\mu,\lambda}\big) \subset \sum_{\gamma \in P} {\mathcal F}_3^{[\mu_a],[\lambda_a],[\gamma_a]},
\]
where by definition $[\mu_a] = (0,\dots,0,\mu,0,\dots,0)$ with $\mu$ on the $a$-th entry, and similarly $[\lambda_a] = (0,\dots,0,\lambda,0,\dots,0)$ and $[\gamma_a] = (0,\dots,0,\gamma,0,\dots,0)$. Therefore, the relation \eqref{echangegrad} yields in ${\rm Gr}_{{\mathcal F}_3}({\mathcal L}_{0,n})$ similar relations between elements of the form (matrix coefficient)$\otimes 1$, or $1 \otimes$(matrix coefficient).

We now consider the case of ``mixed'' products. We give the details when $n=2$, the general case being similar. Let us write the twist $F$ in \eqref{prodL0nform2} as
\[F = \sum_{(F)} F_{(1)}\otimes F_{(2)} = \sum_{(F)} F_{(1)1}\otimes F_{(1)2}\otimes F_{(2)1}\otimes F_{(2)2},\]
that is, we set $F_{(1)1} := R^2_{(2)}R^4_{(2)}$, $F_{(1)2} := R^1_{(2)}R^3_{(2)}$, $F_{(2)1} := R^1_{(1)}R^2_{(1)}$, $F_{(2)2} := R^3_{(1)}R^4_{(1)}$. Put~${d(\mu):= \dim(V_{\mu})}$, $\mu\in P_+$, and
\[\Delta^{(2)}\big({}_{{\mu_2}}\phi^{k_2}_{l_2}\big) = \sum_{p,s=1}^{{ d(\mu_2)}}{}_{{\mu_2}}\phi^{k_2}_{p}\otimes {}_{{\mu_2}}\phi^{p}_{s}\otimes {}_{{\mu_2}}\phi^{s}_{l_2},\qquad \Delta^{(2)}\big({}_{{\mu_1'}}\phi^{k_1'}_{l_1'}\big) = \sum_{p',s'=1}^{{ d(\mu_1')}}{}_{{\mu_1'}}\phi^{k_1'}_{p'}\otimes {}_{{\mu_1'}}\phi^{p'}_{s'}\otimes {}_{{\mu_1'}}\phi^{s'}_{l_1'}.\]
From \eqref{prodL0nform-duplicate}, one obtains
\begin{align}
\big(1 \otimes {}_{{\mu_2}}\phi^{k_2}_{l_2}\big) \big({}_{{\mu_1'}}\phi^{k_1'}_{l_1'} \otimes 1\big)
 ={}& \sum_{(F)} \sum_{p,s=1}^{\scriptscriptstyle{ d(\mu_2)}} \sum_{p',s'=1}^{\scriptscriptstyle{ d(\mu_1')}} \big({}_{{\mu_1'}}\phi^{p'}_{s'} \big({}_{{\mu_1'}}\phi^{k_1'}_{p'}(F_{(2)1}){}_{{\mu_1'}}\phi^{s'}_{l_1'}(S(F_{(2)2}))\big)\big) \nonumber\\
&\otimes \big({}_{{\mu_2}}\phi^{p}_{s} \big({}_{{\mu_2}}\phi^{k_2}_{p}(F_{(1)1}){}_{{\mu_2}}\phi^{s}_{l_2}(S(F_{(1)2}))\big)\big).\label{identcoefnew}
\end{align}
It is immediate that
\[{}_{{\mu_1'}}\phi^{p'}_{s'} \otimes {}_{{\mu_2}}\phi^{p}_{s}\in C(\mu_1')_{w(e_{s'}^{\mu_1'}),w(e_{s'}^{\mu_1'})-w(e_{p'}^{\mu_1'})}\otimes C(\mu_2)_{w(e_{s}^{\mu_2}),w(e_{s}^{\mu_2})-w(e_{p}^{\mu_2})}.\]
As in (iv) above, for every $\mu\in P_+$ we order the weight vectors $e_1^\mu,\dots,e_m^\mu$ so that $w\big(e_i^\mu\big) > w\big(e_j^\mu\big)$ implies $i < j$. With such an ordering the factorization \smash{$R=\Theta\hat{R}$} (see~\eqref{Rmatfact}) implies
\[
{}_{{\mu_2}}\phi^{k_2}_{p}(F_{(1)1}){}_{{\mu_2}}\phi^{s}_{l_2}(S(F_{(1)2}))=0 \qquad \text{unless $k_2\geq p$ and $s\geq l_2$},
\]
and
 \[
 {}_{{\mu_1'}}\phi^{k_1'}_{p'}(F_{(2)1}){}_{{\mu_1'}}\phi^{s'}_{l_1'}(S(F_{(2)2}))=0\qquad \text{unless $k_1'\leq p'$ and $s'\leq l_1'$}.
 \]
Since $s>l_2$, we have $w\big(e_{s}^{\mu_2}\big)\leq w\big(e_{l_2}^{\mu_2}\big)$, and if $w\big(e_{s}^{\mu_2}\big) < w\big(e_{l_2}^{\mu_2}\big)$, then \smash{${}_{{\mu_2}}\phi^{p}_{s} \in \mathcal{F}_2^{< \mu_2,w(e_{l_2}^{\mu_2})}$}. In this last situation, the summands \smash{${}_{{\mu_1'}}\phi^{p'}_{s'} \otimes {}_{{\mu_2}}\phi^{p}_{s}$} in the sum above vanish in ${\rm Gr}_{\mathcal{F}_3}(\Ll_{0,2})$. In order to find all the non-zero summands, we have to consider also the weights with respect to the action~\eqref{coadweightnew}. Since $k_2\geq p$ implies $w\big(e_{k_2}^{\mu_2}\big)\leq w\big(e_{p}^{\mu_2}\big)$, we have~${w\big(e_{s}^{\mu_2}\big)-w\big(e_{p}^{\mu_2}\big) \leq w\big(e_{l_2}^{\mu_2}\big)-w\big(e_{k_2}^{\mu_2}\big)}$. Therefore, the summands which are non-zero in ${{\rm Gr}_{\mathcal{F}_3}(\Ll_{0,2})}$ have both weights $w\big(e_{s}^{\mu_2}\big) = w\big(e_{l_2}^{\mu_2}\big)$ and $w\big(e_{p}^{\mu_2}\big) = w\big(e_{k_2}^{\mu_2}\big)$. Doing similarly with the weights of~\smash{${}_{{\mu_1'}}\phi^{p'}_{s'}$}, we find that also \smash{$w\big(e_{s'}^{\mu_1'}\big) = w\big(e_{l_1'}^{\mu_1'}\big)$} and \smash{$w\big(e_{p'}^{\mu_1'}\big) = w\big(e_{k_1'}^{\mu_1'}\big)$}. When all these conditions on weights are satisfied, the corresponding components $F_{(1)1}$, $F_{(1)2}$, $F_{(2)1}$, $F_{(2)2}$ have zero weight. Therefore, the sum reduces to
\begin{align*}
&\sum_{(F)} {}_{{\mu_2}}\phi_{k_2}^{k_2}(F_{(1)1}){}_{{\mu_2}}\phi_{l_2}^{l_2}(S(F_{(1)2})){}_{{\mu_1'}}\phi_{k_1'}^{k_1'}(F_{(2)1}){}_{{\mu_1'}}\phi_{l_1'}^{l_1'}(S(F_{(2)2}) \\ & \qquad =\big\langle {}_{{\mu_2}}\phi_{k_2}^{k_2}\otimes {}_{{\mu_2}}\phi_{l_2}^{l_2}\otimes {}_{{\mu_1'}}\phi_{k_1'}^{k_1'}\otimes {}_{{\mu_1'}}\phi_{l_1'}^{l_1'}, \Theta_{13}\Theta_{14}^{-1}\Theta_{24}\Theta_{23}^{-1}\big\rangle = q^{(w(e_{k_2}^{\mu_2})-w(e_{l_2}^{\mu_2}), w(e_{k_1'}^{\mu_1'}) - w(e_{l_1'}^{\mu_1'}) )}.
\end{align*}
Denoting by $q'_{k_2l_2k_1'l_1'}$ this scalar, it follows
\begin{align*}
\big(1 \otimes {}_{{\mu_2}}\phi^{k_2}_{l_2}\big) \circ_2 \big({}_{{\mu_1'}}\phi^{k_1'}_{l_1'} \otimes 1\big) & = q'_{k_2l_2k_1'l_1'} \ {}_{{\mu_1'}}\phi^{k_1'}_{l_1'} \otimes {}_{{\mu_2}}\phi^{k_2}_{l_2} = q'_{k_2l_2k_1'l_1'} \big({}_{{\mu_1'}}\phi^{k_1'}_{l_1'} \otimes 1\big)\circ_2 \big(1 \otimes {}_{{\mu_2}}\phi^{k_2}_{l_2}\big).
\end{align*}
This is the relation analogous to \eqref{echangegrad} for mixed products in ${\rm Gr}_{\mathcal{F}_3}(\Ll_{0,2})$.

Recall that in \eqref{echangegrad+} we denoted by $u_1,\dots,u_M$ the ordered list of matrix coefficients ${}_{\varpi_k}\phi_{i}^{j}$. Let us order in a lexicographic way the elements $u_i\otimes u_j$, i.e., as a sequence \smash{$u_1^{(2)},\dots,u_{M^2}^{(2)}$} such that the following condition holds: if \smash{${}_{\varpi_{l'}}\phi_{s'}^{t'} < {}_{\varpi_{k'}}\phi_{i'}^{j'}$}, or \smash{${}_{\varpi_{l'}}\phi_{s'}^{t'} = {}_{\varpi_{k'}}\phi_{i'}^{j'}$} and \smash{${}_{\varpi_l}\phi_{s}^{t} < {}_{\varpi_k}\phi_{i}^{j}$}, then~\smash{$u_a^{(2)}:={}_{\varpi_k}\phi_{i}^{j} \otimes {}_{\varpi_{k'}}\phi_{i'}^{j'}$} and \smash{$u_b^{(2)}:= {}_{\varpi_l}\phi_{s}^{t} \otimes {}_{\varpi_{l'}}\phi_{s'}^{t'}$} satisfy \smash{$u_b^{(2)} < u_a^{(2)}$}. Then, for this ordering the polynormal relations \eqref{echangegrad+} hold true for all \smash{$1\leq u_j^{(2)} < u_i^{(2)}\leq M^2$}. As described after \eqref{echangegrad++}, it follows that ${\rm Gr}_{\mathcal{F}_3}(\Ll_{0,n})$ is Noetherian. The filtration $\mathcal{F}_3$ being well founded, it implies that~$\Ll_{0,n}$ is Noetherian too.

Finally, we consider the $A$-algebra $\Ll_{0,n}^A$, and prove it is Noetherian. We proceed in exactly the same way as for $\Ll_{0,n}$, changing the generators and replacing key arguments of the steps (i)--(iv) by the corresponding results over $A$. Let us describe these modifications step by step.

First, consider the case $n=1$. Recall the $A$-lattices \smash{${}_A \overset{\raisebox{-1.5pt}{\scriptsize$\smallbullet$}}{C} (\lambda)$} (see \eqref{AClambda}), and the decomposition \eqref{OAweightdecomp} of $\Oo_A$ into weight subspaces. In particular, have a decomposition into weight subspaces for the left coregular action,
\[{}_A \overset{\raisebox{-1.5pt}{\scriptsize$\smallbullet$}}{C} (\lambda) = \bigoplus_{\lambda'\in P} {}_A \overset{\raisebox{-1.5pt}{\scriptsize$\smallbullet$}}{C} (\lambda)_{\lambda'}.\]
Define
\[{}_A\mathcal{F}_2^{\mu,\lambda} := \bigoplus_{(\mu',\lambda')\preceq (\mu,\lambda)} {}_A \overset{\raisebox{-1.5pt}{\scriptsize$\smallbullet$}}{C} (\mu')_{\lambda'}.\]
Recall that every $A$-module of matrix coefficients $({\rm End}({}_A V_\mu))^*$, $\mu\in P_+$, is contained in $\Oo_A(\leq \mu)$, and by inverting over $\mc(q)$ the corresponding linear triangular system between basis elements, and using that the order relation $\preceq$ is finer than $\leq$, we obtain an inclusion
\[
 \bigoplus_{\mu'\preceq \mu}\ {}_A \overset{\raisebox{-1.5pt}{\scriptsize$\smallbullet$}}{C} (\mu') \subset \bigoplus_{\mu'\preceq \mu} \ C(\mu')
\]
 (see \eqref{exactsequence}--\eqref{tensorCOA}). It follows that \smash{${}_A\mathcal{F}_2^{\mu,\lambda} = \mathcal{F}_2^{\mu,\lambda} \cap \Oo_A$}, and therefore, like $\mathcal{F}_2$ the family${}_A\mathcal{F}_2 :=\smash{\big\{{}_A\mathcal{F}_2^{\mu,\lambda}\big\}_{(\mu,\lambda)\in \Lambda}}$ is a well-founded filtration of $\Oo_A$. Put \smash{${}_A\mathcal{F}_2^{\prec \mu,\lambda} = \mathcal{F}_2^{\prec \mu,\lambda} \cap \Oo_A$}, and consider the graded $A$-module \smash{${\rm Gr}_{ {}_A\mathcal{F}_2}\big({\mathcal L}_{0,1}^A\big)$} associated to ${}_A\mathcal{F}_2$. By \eqref{OAsurA}--\eqref{OAfacteurs} and the fact that~${\Oo_A=\Ll_{0,1}^A}$ as an $A$-module, we have the $A$-module decomposition
\[{\rm Gr}_{ {}_A\mathcal{F}_2}\big({\mathcal L}_{0,1}^A\big)=\bigoplus_{(\mu,\lambda)\in \Lambda} {}_AC(\mu)_\lambda,\]
where ${}_AC(\mu)_\lambda$ is the submodule of weight $\lambda$ (for the left coregular action) of
\[{}_AC(\mu) := \left({\rm End}({}_A V_\mu\right))^*.\]
Then, we can proceed as before. By step (i), we deduce that ${}_A\mathcal{F}_2$ is an algebra filtration of $\Ll_{0,1}^A$. By Proposition~\ref{teoLuzstigOAscinde}, the $A$-module \smash{${}_A \overset{\raisebox{-1.5pt}{\scriptsize$\smallbullet$}}{C} (\mu_1+\mu_2)$} has multiplicity one in \smash{${}_A \overset{\raisebox{-1.5pt}{\scriptsize$\smallbullet$}}{C} (\mu_1)\otimes {}_A \overset{\raisebox{-1.5pt}{\scriptsize$\smallbullet$}}{C} (\mu_2)$}. In fact, by step (ii), ${}_AC(\mu_1+\mu_2)$ has multiplicity one in ${}_AC(\mu_1) \bigotimes_A {}_AC(\mu_2)$, so exactly in the same way as \eqref{filtration=}, we obtain in ${\rm Gr}_{ {}_A\mathcal{F}_2}\big({\mathcal L}_{0,1}^A\big)$ the equality
\[%\label{filtration=A}
{}_AC(\mu) \circ{}_AC(\nu) = {}_AC(\mu + \nu).
\]
In step (iii), we fixed a basis of each space $C(\mu)$, consisting of a set of matrix coefficients~$\big\{{}_\mu\phi^i_j\big\}$ with respect to dual basis of weight vectors of the modules $V_\mu$ and $V_\mu^*$. In step (iv), the basis elements of $V_\mu$ and $V_\mu^*$ were ordered by means of the weights, and we used the fact that the matrix coefficients in the spaces $C(\varpi_1),\dots,C(\varpi_m)$ form a generating set of the algebra $ {\rm Gr}_{\mathcal{F}_2}({\mathcal L}_{0,1})$. The only property of the matrix coefficients used in the computations was that they are weight vectors for the left coregular action (and later, in the case $n>1$, for the action~\eqref{coadweightnew}).

We can proceed exactly in the same manner by working with the $A$-modules of matrix coefficients ${}_AC(\mu)$. If one wishes to work at the lever of $\Oo_A$, recall that any set of generators of~$\Oo_A$ generates $\Ll_{0,1}^A$ as well (see the proof of Proposition~\ref{L0NAfgfree}). Then, one can replace the basis~\smash{$\big\{{}_\mu\phi^i_j\big\}$} of each space $C(\mu)$ with the canonical basis~\smash{$\dot{\mathbf{B}} [\mu]^*$} of \smash{${}_A \overset{\raisebox{-1.5pt}{\scriptsize$\smallbullet$}}{C} (\mu)$}, and take the generating set of~$\Oo_A$ formed by the elements in \smash{$\dot{\mathbf{B}} [\varpi_i]^*$}, $i=1,\dots,m$ (see Proposition~\ref{genOA} and the comments thereafter). By the integrality properties satisfied by the $R$-matrix and the twists, all the computations in the proof of steps (iii) and (iv) can be done using such basis elements, and eventually take place over $A$ (see \cite[Propositions 4.10 and 6.9]{BR}). Therefore, we obtain a relation like~\eqref{echangegrad+} with coefficients \smash{$\alpha_{ij}^{st}\in A$}. Since $A$ is a Noetherian ring, again this proves~${ {\rm Gr}_{ {}_A\mathcal{F}_2}\big({\mathcal L}_{0,1}^A\big)}$, whence ${\mathcal L}_{0,1}^A$, are Noetherian rings.

This being done, the adaptation of the proof when $n>1$ is immediate. The filtration $\mathcal{F}_3$ has to be replaced with \smash{${}_A\mathcal{F}_3:=\big\{{}_A{\mathcal F}_3^{[\mu],[\lambda],[\gamma]}\big\}_{([\mu],[\lambda],[\gamma])}$}, where \smash{${}_A{\mathcal F}_3^{[\mu],[\lambda],[\gamma]}$} is the $A$-module defined by
\[
{}_A{\mathcal F}_3^{[\mu],[\lambda],[\gamma]}=\bigoplus_{([\mu'],[\lambda'],[\gamma'])\preceq ([\mu],[\lambda],[\gamma])} {}_A \overset{\raisebox{-1.5pt}{\scriptsize$\smallbullet$}}{C} ([\mu'])_{[\lambda'],[\gamma']},
\]
where
\[
{}_A \overset{\raisebox{-1.5pt}{\scriptsize$\smallbullet$}}{C} ([\mu])_{[\lambda],[\gamma]} ={}_A \overset{\raisebox{-1.5pt}{\scriptsize$\smallbullet$}}{C} (\mu_1)_{\lambda_1,\gamma_1}\bigotimes_A \dots \bigotimes_A {}_A \overset{\raisebox{-1.5pt}{\scriptsize$\smallbullet$}}{C} (\mu_n)_{\lambda_n,\gamma_n},
\]
 and \smash{${}_A \overset{\raisebox{-1.5pt}{\scriptsize$\smallbullet$}}{C} (\mu)_{\lambda,\gamma}$} is the subspace of \smash{${}_A \overset{\raisebox{-1.5pt}{\scriptsize$\smallbullet$}}{C} (\mu)_{\lambda}$} of weight $\gamma$ for the action~\eqref{coadweightnew}. Then the proof proceeds in exactly the same way, replacing in~\eqref{identcoefnew} and all subsequent computations the matrix coefficients by the generators of~$\Oo_A$ provided by Proposition~\ref{genOA}. This concludes the proof.
\end{proof}
